\section{Anexos}\label{sec:anexos}

\subsection{Repositorio y estructura del proyecto}
El código fuente completo, los datos, el modelo y este informe están versionados en el repositorio público \url{https://github.com/ForLess01/Accident_NN}. El cuaderno \path{notebooks/modelo.ipynb} documenta el proyecto de forma ejecutable y es compatible con Google Colab.

\begin{lstlisting}[language={}, numbers=none]
Accident_NN/
|-- app/            Interfaz Streamlit (solo lectura del bundle)
|-- data/
|   |-- raw/        Libros oficiales ONSV + manifiesto SHA-256
|   |-- processed/  base_limpia.parquet y casos de demostracion
|   `-- geo/        GeoJSON de departamentos del Peru
|-- docs/           Procedencia y disponibilidad de los datos
|-- models/final/   Bundle canonico: modelo, encoders, calibrador
|-- notebooks/      modelo.ipynb (cuaderno academico integral)
|-- report/         Informe LaTeX, figuras, tablas y PDF final
|-- scripts/        Construccion del informe y gate de release
|-- src/            Pipeline: auditoria, EDA, entrenamiento
`-- tests/          Suite de regresion (pytest)
\end{lstlisting}

\subsection{Configuración canónica}
\begin{lstlisting}
canonical = {
  "version": "canonical-3.0.0",
  "raw_fields": 21,
  "features": 169,
  "hidden_units": [64, 32],
  "dropout": 0.35,
  "l2": 3e-4,
  "learning_rate": 5e-4,
  "batch_size": 64,
  "seed": 42,
  "fixed_refit_epochs": 14,
  "parameters": 12993,
  "raw_threshold": 0.65,
  "calibrated_threshold": 0.15,
  "calibration": "platt"
}
\end{lstlisting}

\subsection{Artefactos de evidencia}
\begin{itemize}\small
\item Procedencia: \path{data/raw/source_manifest.json} y \path{docs/data_provenance.md}.
\item Proxy: \path{report/tables/personas_target_proxy_identity.json}.
\item Curvas/brecha: \path{report/tables/model_learning_curves.csv};\\
      \path{report/tables/model_generalization_gap.json}.
\item Folds rolling: \path{report/tables/temporal_nested_fold_roles.csv};\\
      \path{report/tables/temporal_nested_outer_metrics.csv}.
\item Estabilidad SHAP: \path{report/tables/final_explainability_stability.csv}.
\item Demos: \path{data/processed/demo_cases.csv} y su manifiesto.
\item PDF: \path{report/build_manifest.json}; la frescura depende de hashes, no mtimes.
\end{itemize}

\subsection{Accesibilidad del PDF}
Se añadieron título, autores, asunto y palabras clave. La toolchain pdfLaTeX disponible no genera ni valida PDF/UA etiquetado de manera confiable sin migración de motor/paquetes; no se simula cumplimiento. La accesibilidad interactiva y las alternativas tabulares se verifican en la aplicación.
